\chapter{Termo de Abertura do Projeto}

O Termo de Abertura do Projeto formaliza a existência do projeto e define seus principais objetivos, partes interessadas, requisitos e indicadores de sucesso. Este documento autoriza oficialmente o desenvolvimento do projeto Micromouse Autônomo com Sistema de Telemetria em Tempo Real, permitindo a aplicação dos recursos organizacionais necessários para sua execução no contexto da disciplina Projeto Integrador de Engenharia 1 da Universidade de Brasília.

\section{Dados do projeto}
\begin{description}
    \item [Nome do Projeto:] Micromouse Autônomo com Sistema de Telemetria em Tempo Real
    \item [Data de abertura:] 08/04/2026
    \item [Código:] PI1-MM-2026
    \item [Patrocinador:] Universidade de Brasília
    \item [Gerente:] Giovanni A. Giampauli / 21/1043647 / 211043647@aluno.unb.br / (12) 99638-9028
\end{description}

\section{Objetivos}

O projeto tem como objetivo desenvolver um robô autônomo do tipo micromouse capaz de navegar em um labirinto desconhecido, identificar paredes, mapear o ambiente e encontrar o centro do labirinto de forma totalmente autônoma, seguindo as especificações estabelecidas pela competição Micromouse regulamentada pelo IEEE.

A solução deverá integrar hardware embarcado, algoritmos de navegação, comunicação sem fio e um sistema web de telemetria em tempo real, permitindo o monitoramento remoto da execução do robô.

Os objetivos específicos do projeto são:

\begin{itemize}
    \item Desenvolver um micromouse com dimensões inferiores a 18 cm $\times$ 18 cm;
    
    \item Implementar navegação autônoma utilizando algoritmos de mapeamento e busca de caminho;
    
    \item Garantir precisão mínima de 1 cm nas medições realizadas pelos sensores;
    
    \item Desenvolver um sistema de telemetria com latência inferior a 1 segundo;
    
    \item Permitir o monitoramento remoto da execução do robô por meio de interface web;
    
    \item Garantir autonomia energética suficiente para concluir o labirinto em até 30 minutos;
    
    \item Finalizar o protótipo funcional até o término do semestre letivo 2026.1.
\end{itemize}

Os objetivos foram definidos seguindo a metodologia SMART, sendo específicos, mensuráveis, realistas, acordados entre as áreas do projeto e limitados temporalmente.

\section{Mercado-alvo}

O projeto possui como público-alvo principal estudantes, pesquisadores e instituições de ensino interessados em robótica educacional, sistemas embarcados e automação.

Também fazem parte do mercado-alvo equipes universitárias participantes de competições de robótica, laboratórios acadêmicos e eventos tecnológicos voltados ao desenvolvimento de soluções autônomas.

Além disso, o projeto pode ser utilizado como plataforma didática para o ensino de programação, eletrônica, controle e inteligência artificial aplicada à robótica móvel.

\section{Requisitos}

Os principais requisitos do projeto incluem:

\begin{itemize}
    \item Capacidade de movimentação autônoma em labirintos;
    
    \item Reconhecimento de paredes e obstáculos utilizando sensores;
    
    \item Precisão mínima de 1 cm na medição de distâncias;
    
    \item Controle dos motores para deslocamento e rotação;
    
    \item Comunicação sem fio entre o robô e o sistema web;
    
    \item Exibição de dados de telemetria em tempo real;
    
    \item Armazenamento do mapa descoberto e da trajetória percorrida;
    
    \item Dimensões inferiores a 18 cm $\times$ 18 cm;
    
    \item Funcionamento contínuo sem reinicializações causadas por oscilações elétricas;
    
    \item Autonomia energética suficiente para completar o desafio proposto.
\end{itemize}

\section{Justificativa}

O desenvolvimento de sistemas robóticos autônomos tem se tornado cada vez mais relevante devido ao crescimento das áreas de automação, inteligência artificial e Internet das Coisas (IoT). Nesse contexto, o micromouse representa um problema clássico de robótica móvel, envolvendo navegação autônoma, tomada de decisão em tempo real e integração entre hardware e software.

Além da relevância tecnológica, o projeto possui forte caráter educacional, permitindo a aplicação prática de conceitos de eletrônica, programação, controle, algoritmos e engenharia de software em um sistema real e integrado.

O projeto também busca democratizar o acesso à robótica competitiva no Brasil por meio do desenvolvimento de uma solução de baixo custo baseada em componentes amplamente disponíveis no mercado nacional.

\section{Indicadores}

Os indicadores definidos para avaliação do projeto são:

\begin{enumerate}
    \item Tempo total para conclusão do labirinto;
    
    \item Quantidade de colisões durante a navegação;
    
    \item Precisão das medições dos sensores;
    
    \item Latência da telemetria em tempo real;
    
    \item Frequência de atualização dos dados enviados;
    
    \item Consumo energético total do sistema;
    
    \item Taxa de sucesso na chegada ao objetivo;
    
    \item Estabilidade elétrica durante a execução;
    
    \item Quantidade de estudantes impactados pelo projeto;
    
    \item Custo total de desenvolvimento comparado a soluções comerciais.
\end{enumerate}